\documentclass[11pt]{article}
\usepackage[margin=0.75in]{geometry}
\usepackage{amsmath,amssymb,booktabs,longtable,array}
\usepackage{xcolor,listings,fancyhdr,hyperref}
\usepackage{enumitem}
\definecolor{navy}{RGB}{24,55,95}
\definecolor{green}{RGB}{25,100,70}
\definecolor{codebg}{RGB}{246,248,250}
\definecolor{codecomment}{RGB}{90,110,95}
\hypersetup{colorlinks=true,linkcolor=navy,urlcolor=navy}
\pagestyle{fancy}
\fancyhf{}
\lhead{STAT452}
\rhead{Comprehensive R Syntax Reference}
\cfoot{\thepage}
\setlength{\headheight}{14pt}
\lstdefinelanguage{Rcustom}{
  language=R,
  morekeywords={lm,glm,glmnet,lm.ridge,summary,coef,predict,confint,
    fitted,resid,residuals,model.matrix,anova,AIC,BIC,step,cv.glm,
    scale,poly,plot,par,table,library,factor,subset},
  sensitive=true
}
\lstset{
  language=Rcustom,
  basicstyle=\ttfamily\small,
  keywordstyle=\color{navy}\bfseries,
  commentstyle=\color{codecomment}\itshape,
  stringstyle=\color{green},
  backgroundcolor=\color{codebg},
  frame=single,
  rulecolor=\color{black!15},
  breaklines=true,
  columns=fullflexible,
  keepspaces=true,
  showstringspaces=false,
  xleftmargin=0.4em,
  xrightmargin=0.4em,
  aboveskip=0.5em,
  belowskip=0.7em
}
\title{STAT452 Comprehensive R Syntax Reference\\
\large Multiple Linear Regression, Polynomial Regression, Regularization,\\
Variable Selection, Logistic Regression, and Classification}
\author{STAT452 -- Applied Statistics for Engineers and Scientists II}
\date{}
\begin{document}
\maketitle
\tableofcontents
\bigskip

\section{How to use this reference}
This is a practical syntax reference for the four course topics. The examples
use built-in \texttt{mtcars} data whenever possible. Run commands one section
at a time in R or RStudio. Lines beginning with \texttt{\#} are comments.

\section{Basic R syntax}

\subsection{Values, assignment, and help}
\begin{lstlisting}
# This is a comment
x <- 5
x == 5
c(1, 2, 3)
seq(0, 10, by = 2)
rep(1:3, times = 2)
length(x)
class(x)
typeof(x)
is.na(x)
?mean
help(lm)
getwd()
setwd("path/to/folder")
set.seed(452)
\end{lstlisting}

Use \texttt{<-} for assignment and \texttt{==} for an equality test. A
function is called with parentheses, for example \texttt{mean(x)}.

\subsection{Vectors, matrices, and data frames}
\begin{lstlisting}
v <- c(2, 4, 6, 8)
v[1]
v[2:4]
v[-1]
v[v > 4]

M <- matrix(1:6, nrow = 2, byrow = TRUE)
M[1, 2]
t(M)
dim(M)
nrow(M)
ncol(M)

d <- data.frame(
  y = c(10, 12, 15),
  x = c(1, 2, 3),
  group = factor(c("A", "A", "B"))
)
d$y
d[, c("y", "x")]
d[1:2, ]
subset(d, x > 1)
d$new_variable <- d$x^2
names(d)
str(d)
head(d)
summary(d)
\end{lstlisting}

\subsection{Data input and preparation}
\begin{lstlisting}
data(mtcars)
cars <- mtcars

d <- read.csv("data.csv")
write.csv(cars, "cars.csv", row.names = FALSE)

na.omit(d)
complete.cases(d)
d[complete.cases(d), ]

factor(c("No", "Yes"), levels = c("No", "Yes"))
factor(x, labels = c("automatic", "manual"))
droplevels(d)

mean(x)
mean(x, na.rm = TRUE)
median(x, na.rm = TRUE)
sd(x, na.rm = TRUE)
var(x, na.rm = TRUE)
min(x)
max(x)
range(x)
quantile(x, probs = c(.25, .5, .75), na.rm = TRUE)
sum(x)
table(d$group)
prop.table(table(d$group))
scale(x)
\end{lstlisting}

\section{Graphics and exploratory analysis}

\begin{lstlisting}
plot(x, y)
plot(y ~ x, data = d)
hist(x, breaks = 10)
boxplot(y ~ group, data = d)
pairs(d[, c("y", "x1", "x2")])
cor(d[, c("x1", "x2", "x3")], use = "complete.obs")

plot(x, y, pch = 19, col = "navy",
     main = "Response versus predictor",
     xlab = "Predictor", ylab = "Response")
abline(lm(y ~ x, data = d), col = "red", lwd = 2)

par(mfrow = c(2, 2))
plot(fit)
par(mfrow = c(1, 1))
\end{lstlisting}

\section{Topic 1: Multiple linear regression}

\subsection{Model formulas}
\begin{lstlisting}
y ~ x
y ~ x1 + x2
y ~ .
y ~ x - 1
y ~ 0 + x
y ~ x1 * x2
y ~ x1:x2
y ~ x + factor(group)
\end{lstlisting}

The formula \texttt{y \textasciitilde{} x1 + x2} means an intercept plus
two predictors. The formula \texttt{x1 * x2} includes both main effects and
their interaction. The formula \texttt{x1:x2} includes the interaction only.

\subsection{Fit, summarize, and inspect}
\begin{lstlisting}
fit <- lm(mpg ~ wt + hp, data = mtcars)
fit
summary(fit)
coef(fit)
coefficients(fit)
confint(fit)
vcov(fit)
formula(fit)
model.frame(fit)
model.matrix(fit)
\end{lstlisting}

\subsection{Coefficients and fitted values}
\begin{lstlisting}
coef(fit)["wt"]
coef(summary(fit))
summary(fit)$coefficients
summary(fit)$sigma
summary(fit)$r.squared
summary(fit)$adj.r.squared
df.residual(fit)

y_hat <- fitted(fit)
e <- resid(fit)
head(data.frame(
  observed = mtcars$mpg,
  fitted = y_hat,
  residual = e
))
max(abs(mtcars$mpg - (y_hat + e)))
sum(e)
\end{lstlisting}

For a multiple-regression coefficient, always state that the other predictors
are held fixed.

\subsection{Manual simple-regression calculations}
\begin{lstlisting}
x <- mtcars$wt
y <- mtcars$mpg
x_bar <- mean(x)
y_bar <- mean(y)

Sxx <- sum((x - x_bar)^2)
Sxy <- sum((x - x_bar) * (y - y_bar))
b1 <- Sxy / Sxx
b0 <- y_bar - b1 * x_bar
c(b0 = b0, b1 = b1)
\end{lstlisting}

\subsection{SSE, MSE, and $R^2$}
\begin{lstlisting}
SSE <- sum(e^2)
SST <- sum((y - mean(y))^2)
SSR <- SST - SSE
MSE <- SSE / df.residual(fit)
RMSE <- sqrt(MSE)
R2 <- 1 - SSE / SST

c(SSE = SSE, SST = SST, SSR = SSR,
  MSE = MSE, RMSE = RMSE, R2 = R2)
\end{lstlisting}

\subsection{Matrix least squares}
The least-squares estimator is
\[
 \widehat{\boldsymbol\beta}
 =(\mathbf X^{\mathsf T}\mathbf X)^{-1}
   \mathbf X^{\mathsf T}\mathbf Y.
\]
\begin{lstlisting}
X <- model.matrix(fit)
Y <- matrix(mtcars$mpg, ncol = 1)

t(X)
t(X) %*% X
solve(t(X) %*% X)
beta_hat <- solve(t(X) %*% X, t(X) %*% Y)
beta_hat
coef(fit)

drop(t(X) %*% matrix(resid(fit), ncol = 1))
\end{lstlisting}

Use \texttt{\%*\%} for matrix multiplication. Use \texttt{*} for
element-by-element multiplication.

\subsection{Coefficient tests and confidence intervals}
\begin{lstlisting}
tab <- coef(summary(fit))
estimate <- tab["hp", "Estimate"]
standard_error <- tab["hp", "Std. Error"]
t_value <- estimate / standard_error
df_error <- df.residual(fit)
p_value <- 2 * pt(abs(t_value),
                  df = df_error,
                  lower.tail = FALSE)
c(t = t_value, df = df_error, p = p_value)
confint(fit, "hp", level = 0.95)
\end{lstlisting}

\subsection{Prediction}
\begin{lstlisting}
new_car <- data.frame(wt = 3, hp = 110)

predict(fit, newdata = new_car)
predict(fit, newdata = new_car,
        interval = "confidence", level = 0.95)
predict(fit, newdata = new_car,
        interval = "prediction", level = 0.95)
\end{lstlisting}

A confidence interval estimates the mean response. A prediction interval
estimates one new individual response and is wider.

\subsection{Nested models and partial F-tests}
\begin{lstlisting}
reduced <- lm(mpg ~ wt, data = mtcars)
full <- lm(mpg ~ wt + hp, data = mtcars)

anova(reduced, full)
anova(full)
drop1(full, test = "F")
\end{lstlisting}

\subsection{Diagnostics and influence}
\begin{lstlisting}
par(mfrow = c(2, 2))
plot(fit)
par(mfrow = c(1, 1))

rstandard(fit)
rstudent(fit)
hatvalues(fit)
cooks.distance(fit)
influence.measures(fit)
dfbetas(fit)

which.max(abs(rstandard(fit)))
which.max(cooks.distance(fit))
\end{lstlisting}

The four standard diagnostic plots check linearity and spread, normality,
constant variance, and leverage/influence.

\section{Topic 2: Polynomial regression and transformations}

\subsection{Polynomial terms}
\begin{lstlisting}
linear <- lm(mpg ~ wt, data = mtcars)
quadratic <- lm(mpg ~ wt + I(wt^2), data = mtcars)
cubic <- lm(mpg ~ wt + I(wt^2) + I(wt^3), data = mtcars)

summary(quadratic)
anova(linear, quadratic, cubic)
AIC(linear, quadratic, cubic)
BIC(linear, quadratic, cubic)
\end{lstlisting}

Use \texttt{I(wt\textasciicircum2)} to make R evaluate the arithmetic
expression inside a formula. Retain lower-order terms when higher-order terms
are included.

\subsection{Orthogonal polynomials}
\begin{lstlisting}
orthogonal_fit <- lm(mpg ~ poly(wt, degree = 2), data = mtcars)
raw_fit <- lm(mpg ~ poly(wt, degree = 2, raw = TRUE),
              data = mtcars)
summary(orthogonal_fit)
summary(raw_fit)
\end{lstlisting}

\texttt{poly(wt, 2)} uses an orthogonal polynomial basis by default.
\texttt{raw = TRUE} gives ordinary powers of \texttt{wt}.

\subsection{Centering and stationary points}
\begin{lstlisting}
mtcars$wt_c <- mtcars$wt - mean(mtcars$wt)
centered_fit <- lm(mpg ~ wt_c + I(wt_c^2), data = mtcars)

b <- coef(centered_fit)
turning_point_centered <-
  -b["wt_c"] / (2 * b["I(wt_c^2)"])
turning_point_original <-
  turning_point_centered + mean(mtcars$wt)
\end{lstlisting}

\subsection{Polynomial predictions}
\begin{lstlisting}
grid <- data.frame(
  wt = seq(min(mtcars$wt),
           max(mtcars$wt),
           length.out = 100)
)
grid$fit <- predict(quadratic, newdata = grid)

plot(mpg ~ wt, data = mtcars, pch = 19)
lines(grid$wt, grid$fit, col = "blue", lwd = 2)

predict(quadratic, newdata = data.frame(wt = c(1, 3, 6)))
range(mtcars$wt)
\end{lstlisting}

Predictions outside the observed range are extrapolations and may be highly
unstable.

\subsection{Transformations}
\begin{lstlisting}
log_y_fit <- lm(log(mpg) ~ wt, data = mtcars)
log_x_fit <- lm(mpg ~ log(wt), data = mtcars)
log_log_fit <- lm(log(mpg) ~ log(wt), data = mtcars)

sqrt_fit <- lm(sqrt(mpg) ~ wt, data = mtcars)
inverse_fit <- lm(mpg ~ I(1 / wt), data = mtcars)
power_fit <- lm(mpg ~ I(wt^0.5), data = mtcars)

summary(log_log_fit)
par(mfrow = c(2, 2))
plot(log_log_fit)
par(mfrow = c(1, 1))
\end{lstlisting}

Common transformations are \texttt{log(x)}, \texttt{sqrt(x)}, \texttt{x\^2},
\texttt{1/x}, and \texttt{exp(x)}. Check residuals after transforming.

\section{Topic 3: Ridge, LASSO, and variable selection}

\subsection{OLS baseline and correlation}
\begin{lstlisting}
full <- lm(mpg ~ wt + hp + qsec + disp + drat + gear,
           data = mtcars)
summary(full)

predictors <- mtcars[, c("wt", "hp", "qsec",
                         "disp", "drat", "gear")]
cor(predictors)
X_standardized <- scale(predictors)
\end{lstlisting}

\subsection{Ridge regression with \texttt{MASS}}
\begin{lstlisting}
library(MASS)

lambda_grid <- seq(0, 30, length.out = 100)
ridge_fit <- lm.ridge(
  mpg ~ wt + hp + qsec + disp + drat + gear,
  data = mtcars,
  lambda = lambda_grid
)

ridge_fit$coef
ridge_fit$GCV
best_index <- which.min(ridge_fit$GCV)
lambda_grid[best_index]
coef(ridge_fit)[best_index, ]
plot(ridge_fit)
select(ridge_fit)
\end{lstlisting}

Ridge uses an L2 penalty. Coefficients shrink smoothly toward zero but
normally do not become exactly zero.

\subsection{Ridge and LASSO with \texttt{glmnet}}
If the \texttt{glmnet} package is installed, use:
\begin{lstlisting}
library(glmnet)

x <- model.matrix(
  mpg ~ wt + hp + qsec + disp + drat + gear,
  data = mtcars
)[, -1]
y <- mtcars$mpg

ridge_glmnet <- glmnet(x, y, alpha = 0)
lasso_glmnet <- glmnet(x, y, alpha = 1)

plot(ridge_glmnet, xvar = "lambda")
plot(lasso_glmnet, xvar = "lambda")

cv_ridge <- cv.glmnet(x, y, alpha = 0, nfolds = 5)
cv_lasso <- cv.glmnet(x, y, alpha = 1, nfolds = 5)

cv_lasso$lambda.min
cv_lasso$lambda.1se
coef(cv_lasso, s = "lambda.min")
predict(cv_lasso, newx = x, s = "lambda.min")
\end{lstlisting}

LASSO uses an L1 penalty and can set coefficients exactly to zero. The
\texttt{alpha} argument is \texttt{0} for ridge and \texttt{1} for LASSO.

\subsection{Soft thresholding}
\begin{lstlisting}
soft_threshold <- function(z, gamma) {
  sign(z) * max(abs(z) - gamma, 0)
}
soft_threshold(2.5, 1)
soft_threshold(-2.5, 1)
\end{lstlisting}

Soft thresholding is the basic coordinate-descent update behind a simple
LASSO implementation.

\subsection{AIC, BIC, and stepwise selection}
\begin{lstlisting}
reduced <- lm(mpg ~ wt + hp, data = mtcars)

AIC(reduced, full)
BIC(reduced, full)
anova(reduced, full)

step(full, direction = "forward", trace = 0)
step(full, direction = "backward", trace = 0)
step(full, direction = "both", trace = 0)

bic_model <- step(
  full,
  direction = "both",
  k = log(nrow(mtcars)),
  trace = 0
)
formula(bic_model)
\end{lstlisting}

AIC generally uses a lighter complexity penalty. BIC uses
\texttt{k = log(n)} and is usually more parsimonious.

\subsection{Cross-validation}
\begin{lstlisting}
set.seed(452)
fold <- sample(rep(1:5, length.out = nrow(mtcars)))

cv_mse <- function(formula, data, fold) {
  mean(sapply(sort(unique(fold)), function(k) {
    train <- data[fold != k, ]
    test <- data[fold == k, ]
    fit <- lm(formula, data = train)
    mean((test$mpg -
          predict(fit, newdata = test))^2)
  }))
}

cv_mse(mpg ~ wt, mtcars, fold)
cv_mse(mpg ~ wt + hp + qsec + disp + drat + gear,
       mtcars, fold)
\end{lstlisting}

Cross-validation estimates performance on observations not used to fit the
corresponding model. It is more informative for prediction than training SSE.

\subsection{Leave-one-out cross-validation}
\begin{lstlisting}
library(boot)

loocv <- cv.glm(
  data = mtcars,
  glmfit = glm(mpg ~ wt + hp, data = mtcars),
  K = nrow(mtcars)
)
loocv$delta
\end{lstlisting}

\section{Topic 4: Logistic regression and classification}

\subsection{Create a binary response}
\begin{lstlisting}
data(mtcars)

cars <- transform(
  mtcars,
  transmission = factor(
    am,
    levels = c(0, 1),
    labels = c("automatic", "manual")
  )
)
table(cars$transmission)
\end{lstlisting}

\subsection{Fit a logistic model}
\begin{lstlisting}
logit_fit <- glm(
  transmission ~ wt + hp,
  data = cars,
  family = binomial(link = "logit")
)

summary(logit_fit)
coef(logit_fit)
confint.default(logit_fit)
\end{lstlisting}

The logistic model is
\[
 \log\left(\frac{\pi(x)}{1-\pi(x)}\right)
 =\beta_0+\beta_1x_1+\cdots+\beta_px_p,
\]
where $\pi(x)=P(Y=1\mid x)$.

\subsection{Odds ratios}
\begin{lstlisting}
odds_ratios <- exp(coef(logit_fit))
odds_ratios

odds_ratio_intervals <- exp(confint.default(logit_fit))
odds_ratio_intervals

exp(10 * coef(logit_fit)["hp"])
\end{lstlisting}

\texttt{exp(beta)} is the multiplicative change in odds for a one-unit
increase in a predictor, holding the other predictors fixed.

\subsection{Predicted probabilities}
\begin{lstlisting}
prob_manual <- predict(logit_fit, type = "response")
linear_predictor <- predict(logit_fit, type = "link")

head(prob_manual)
range(prob_manual)

new_car <- data.frame(wt = 3, hp = 110)
predict(logit_fit, newdata = new_car,
        type = "response")
predict(logit_fit, newdata = new_car,
        type = "link")
\end{lstlisting}

Use \texttt{type = "response"} for probabilities and
\texttt{type = "link"} for fitted log-odds.

\subsection{Classification}
\begin{lstlisting}
threshold <- 0.5
predicted_class <- ifelse(
  prob_manual >= threshold,
  "manual",
  "automatic"
)
predicted_class <- factor(
  predicted_class,
  levels = levels(cars$transmission)
)

table(
  observed = cars$transmission,
  predicted = predicted_class
)
mean(predicted_class == cars$transmission)
\end{lstlisting}

\subsection{Confusion-matrix metrics}
\begin{lstlisting}
confusion <- table(
  observed = cars$transmission,
  predicted = predicted_class
)

TP <- confusion["manual", "manual"]
TN <- confusion["automatic", "automatic"]
FP <- confusion["automatic", "manual"]
FN <- confusion["manual", "automatic"]

accuracy <- (TP + TN) / sum(confusion)
sensitivity <- TP / (TP + FN)
specificity <- TN / (TN + FP)
precision <- TP / (TP + FP)

c(accuracy = accuracy,
  sensitivity = sensitivity,
  specificity = specificity,
  precision = precision)
\end{lstlisting}

Sensitivity is the true-positive rate. Specificity is the true-negative rate.
The best threshold depends on the relative costs of false positives and false
negatives.

\subsection{Compare thresholds}
\begin{lstlisting}
classify_at <- function(probability, truth, threshold = 0.5) {
  prediction <- factor(
    ifelse(probability >= threshold,
           "manual", "automatic"),
    levels = levels(truth)
  )
  table(observed = truth, predicted = prediction)
}

classify_at(prob_manual, cars$transmission, 0.3)
classify_at(prob_manual, cars$transmission, 0.5)
classify_at(prob_manual, cars$transmission, 0.7)
\end{lstlisting}

\subsection{Deviance, AIC, and nested logistic models}
\begin{lstlisting}
null_fit <- glm(
  transmission ~ 1,
  data = cars,
  family = binomial
)

anova(null_fit, logit_fit, test = "Chisq")
AIC(null_fit, logit_fit)
BIC(null_fit, logit_fit)
deviance(logit_fit)
logLik(logit_fit)
\end{lstlisting}

\subsection{Logistic diagnostics}
\begin{lstlisting}
plot(logit_fit)
residuals(logit_fit, type = "deviance")
residuals(logit_fit, type = "pearson")
hatvalues(logit_fit)
cooks.distance(logit_fit)
\end{lstlisting}

Check for fitted probabilities near zero or one, complete or quasi-complete
separation, influential observations, and poor calibration.

\section{Training/test validation}

\begin{lstlisting}
set.seed(452)
test_index <- sample(
  seq_len(nrow(cars)),
  size = 8
)

train <- cars[-test_index, ]
test <- cars[test_index, ]

train_fit <- glm(
  transmission ~ wt + hp,
  data = train,
  family = binomial
)

test_probability <- predict(
  train_fit,
  newdata = test,
  type = "response"
)

test_prediction <- ifelse(
  test_probability >= 0.5,
  "manual",
  "automatic"
)

mean(test_prediction == test$transmission)
table(
  observed = test$transmission,
  predicted = test_prediction
)
\end{lstlisting}

Training accuracy is often optimistic because the model was fitted on those
observations. Test-set or cross-validation performance is more informative
about generalization.

\section{Common mistakes and final checklist}

\begin{lstlisting}
# Correct model syntax
fit <- lm(y ~ x1 + x2, data = d)

# Polynomial term
fit2 <- lm(y ~ x + I(x^2), data = d)

# Matrix multiplication
t(X) %*% X

# Logistic probability
predict(logit_fit, type = "response")

# New data must use original predictor names
new_data <- data.frame(x1 = 3, x2 = 110)
predict(fit, newdata = new_data)
\end{lstlisting}

\begin{itemize}[leftmargin=*]
\item Use \texttt{\textasciitilde{}} to define a model formula.
\item Use \texttt{I(x\textasciicircum2)} or \texttt{poly(x, 2)} for polynomial terms.
\item Use \texttt{\%*\%}, not \texttt{*}, for matrix multiplication.
\item Standardize predictors before ridge or LASSO.
\item Use \texttt{type = "response"} for logistic probabilities.
\item Match factor levels and predictor names between training and new data.
\item Interpret multiple-regression slopes while holding other predictors fixed.
\item Do not use p-values as the only variable-selection criterion.
\item Do not delete influential observations automatically.
\item Do not extrapolate polynomial models casually.
\end{itemize}

\[
\boxed{\text{inspect} \longrightarrow \text{fit} \longrightarrow
\text{diagnose} \longrightarrow \text{validate} \longrightarrow
\text{interpret}}
\]

\end{document}

